\section{Where do bad curves come from?}\label{sec-diagnostic}

\begin{objectif}
(G1) Decompose the implied repo, compute the size of the normal noise, and give the shape that each mechanism causing
real anomalies leaves in the carry difference.
\end{objectif}

\subsection{The panel gap}\label{sec-gap}
We work under \cref{hyp-coherent,hyp-panels,hyp-consensus} at a fixed maturity $T$. The contributors' views are
perturbations of a common centre, $q_j=q+u_j$ and $D_j=D+v_j$, where $u_j$ and $v_j$ are of order $\epsilon$. Let
$\qbar$ and $\Dopt$ be the averages of the $q_j$ and of the $D_j$ over the option panel, and $\Ddiv$ the consensus of
the dividend panel, allocated at maturity $T$.

\begin{proposition}[Decomposition of the implied repo]\label{prop-gap}
\begin{equation}\label{eq-gap}
q^{\tot}(T)=\qbar(T)+\frac{\Dopt(T)-\Ddiv(T)}{T\,\bigl(S_0-D(T)\bigr)}+O(\epsilon^2).
\end{equation}
\end{proposition}
\begin{proof}
For $j\in\Popt$, \cref{lem-logf} gives $\ln F_j=\ln F-T\,u_j-v_j/(S_0-D)+O(\epsilon^2)$. The Totem forward is the
average of the $F_j$, and the log of an average equals the average of the logs up to a term of order $\epsilon^2$; hence
$\ln F^{\tot}=\ln F-T\,\bar u^{\opt}-\bar v^{\opt}/(S_0-D)+O(\epsilon^2)$. The controller takes out the consensus
dividends $D^{\tot}=D+\bar v^{\dv}$:
\[
q^{\tot}T=rT-\ln F^{\tot}+\ln\bigl(S_0-D^{\tot}\bigr)=qT+T\,\bar u^{\opt}+\frac{\bar v^{\opt}-\bar v^{\dv}}{S_0-D}
+O(\epsilon^2).
\]
Divide by $T$, with $q+\bar u^{\opt}=\qbar$ and $\bar v^{\opt}-\bar v^{\dv}=\Dopt-\Ddiv$.
\end{proof}

The implied repo is therefore the average repo of the option panel plus a \emph{panel gap}: the difference between the
option panel's average dividend view, which nobody observes, and the dividend panel's, divided by $T(S_0-D)$. For a
given contributor, an error on its dividends cancels exactly in its implied repo; between two panels, it does not.

\begin{corollary}[Size of the panel gap]\label{cor-gap}
Assume the view errors on $D(T)$ are independent and identically distributed with standard deviation $\sigma_D(T)$ in
both panels, and the $n^{\com}$ contributors common to both services consistent with themselves. Then
\begin{equation}\label{eq-gap-sd}
\sigma\bigl(q^{\tot}-\qbar\bigr)=\frac{\sigma_D(T)}{T\,\bigl(S_0-D(T)\bigr)}
\sqrt{\frac{1}{n^{\opt}}+\frac{1}{n^{\dv}}-\frac{2\,n^{\com}}{n^{\opt}n^{\dv}}} .
\end{equation}
\end{corollary}
\begin{proof}
$\Dopt-\Ddiv=\sum_j(\alpha_j-\gamma_j)\,v_j$, with $\alpha_j=1/n^{\opt}$ on the option panel, $\gamma_j=1/n^{\dv}$ on the
dividend panel, and $0$ elsewhere. Its variance is
$\sigma_D^2\sum_j(\alpha_j-\gamma_j)^2=\sigma_D^2\bigl[n^{\com}(1/n^{\opt}-1/n^{\dv})^2+(n^{\opt}-n^{\com})/(n^{\opt})^2
+(n^{\dv}-n^{\com})/(n^{\dv})^2\bigr]$, which simplifies to
$\sigma_D^2\bigl(1/n^{\opt}+1/n^{\dv}-2n^{\com}/(n^{\opt}n^{\dv})\bigr)$.
\end{proof}

The gap vanishes if the two panels are identical ($n^{\com}=n^{\opt}=n^{\dv}$); it is largest when they are disjoint. The
dispersion $\sigma_D(T)$ of a cumulative dividend grows at least like $T$: the panel gap is small in carry at the short
end, but it does not shrink in repo at the long end.

\begin{example}[Size of the gap for a stock]\label{ex-gap}
We take the running example's model: spot $100$, an annual dividend of $\num{4.2}$ going ex in early May, and each
contributor wrong on the dividend going ex at $t$ by a relative error with standard deviation $\num{0.01}+\num{0.03}\,t$,
part of it common to all its dividends. The panels are disjoint, with the running example's sizes, and the option
contributors also disagree by $15$\,bp on the repo, which the implied repo must not pick up. Formula
\eqref{eq-gap-sd} and a contributor-level simulation give:
\begin{center}\input{generated/\lang/tab_gap}\end{center}
The gap alone produces \vGapOneY\,bp of repo noise at one year and \vGapFiveY\,bp at five years, with no error at all
in Totem. It is zero before the first ex-date, jumps at each ex-date, then decays like $1/T$ until the next one
(\cref{fig-gap}).
\end{example}

\begin{figure}[ht]
\centering\includegraphics[width=\linewidth]{figures/\lang/panel_gap.pdf}
\caption{(a) Implied repo noise from the panel gap alone, formula \eqref{eq-gap-sd}, for the running example's panels
and two other sizes (options / dividends): it appears at each ex-date, then decays like $1/T$. (b) A spot mismatch
$\varepsilon$ is constant in carry, hence $\varepsilon/T$ in repo (\cref{prop-signatures}).}
\label{fig-gap}
\end{figure}

\subsection{Contributor drop-offs}\label{sec-nested}
Panels shrink with maturity. Since they are nested, the average over a small panel is not independent of the average
over a larger panel that contains it.

\begin{proposition}[Averages over nested panels]\label{prop-nested}
Let $e_{jk}$ be the error of contributor $j$ at maturity $T_k$, independent across contributors, with
$\Var e_{jk}=\sigma_k^2$ and $\Cov(e_{jk}, e_{jl})=\rho_{kl}\,\sigma_k\sigma_l$. If the panels are nested, with sizes
$n_k$, the averages $\bar e_k$ satisfy
\begin{equation}\label{eq-nested}
\Cov(\bar e_k, \bar e_l)=\frac{\rho_{kl}\,\sigma_k\sigma_l}{\max(n_k, n_l)} .
\end{equation}
In particular, for persistent views ($\rho_{kl}=1$, $\sigma_k=\sigma_l=\sigma$) and $n_l\le n_k$, the jump of the
average when the panel goes from $n_k$ to $n_l$ contributors has standard deviation $\sigma\sqrt{1/n_l-1/n_k}$.
\end{proposition}
\begin{proof}
The two averages share $\min(n_k, n_l)$ contributors, hence
\[
\Cov(\bar e_k, \bar e_l)=\frac{\min(n_k, n_l)}{n_kn_l}\,\rho_{kl}\sigma_k\sigma_l
=\frac{\rho_{kl}\,\sigma_k\sigma_l}{\max(n_k, n_l)} .
\]
Then $\Var(\bar e_l-\bar e_k)=\sigma^2(1/n_l+1/n_k-2/n_k)=\sigma^2(1/n_l-1/n_k)$.
\end{proof}

In units of contributor dispersion, the consensus thus jumps by itself by $\vJumpA\,\sigma$ when the panel goes from
$16$ to $12$ contributors, by $\vJumpB\,\sigma$ from $12$ to $7$ (the running example's option panel between two and
three years), by $\vJumpC\,\sigma$ from $15$ to $7$ and by $\vJumpD\,\sigma$ from $10$ to $4$. A jump of that size at a
contributor drop-off is noise, not information; the noise model of \cref{sec-detection} reproduces it through
\eqref{eq-nested}.

\subsection{Six mechanisms, six signatures}\label{sec-signatures}
Beyond the noise, some mechanisms cause real anomalies. Each leaves a characteristic shape in the carry difference.

\begin{proposition}[Signatures]\label{prop-signatures}
To first order, each mechanism adds the following contribution to the carry difference $\delta$.
\begin{enumerate}[label=(\roman*)]
\item \emph{Spot mismatch.} If the option contributors price off a spot $S_0(1+\varepsilon)$ while the controller uses
$S_0$: $-\varepsilon S_0/(S_0-D_k)\approx-\varepsilon$ at every pillar, that is, a repo error $-\varepsilon/T_k$. A
dividend going ex before $T_1$ and wrong by $\delta d$ in the dividend service gives the same shape,
$-\delta\tilde d/(S_0-D_k)$.
\item \emph{Error on one dividend period.} If the Totem level of period $p$ is wrong by $\delta\ell_p$, the
contribution is $-a_{kp}\,\delta\ell_p/(S_0-D_k)$: a step that starts at the first pillar after the ex-date; in repo, a
jump that decays like $1/T$.
\item \emph{Ex-date disagreement.} Take a dividend going ex at $t_i\in\intoo{T_k}{T_{k+1}}$, and a fraction $f$ of the
option panel that puts it before $T_k$: $+f\,\tilde d_i/(S_0-D_k)$ at pillar $T_k$ alone, a spike. The variance of the
$\ln F_j$ at $T_k$ rises by $f(1-f)\,\tilde d_i^{\,2}/(S_0-D_k)^2$.
\item \emph{Stale ex-date profile.} If the forwards use the true dates $t_i$ but the allocation puts them at
$t_i+\Delta$ with $\Delta>0$: $+\tilde d_i/(S_0-D_k)$ at the pillars $T_k\in\intco{t_i}{t_i+\Delta}$, and nothing
elsewhere (for $\Delta<0$, $-\tilde d_i/(S_0-D_k)$ at the pillars $T_k\in\intco{t_i+\Delta}{t_i}$). For a seasonal
payer, the affected pillars come back every year.
\item \emph{Forward outlier.} If the consensus $\ln F^{\tot}_k$ is off by $\mu$: $-\mu$ at pillar $T_k$ alone, a
spike.
\item \emph{Proportional dividend bias} (for example gross versus net of withholding tax). If every Totem dividend is
$(1+\tau)$ times the one in the forwards: $-\tau D_k/(S_0-D_k)$, a staircase that steps up at every ex-date. For a
continuous stream with yield $y$, $D_k\approx S_0\,y\,T_k$ and the contribution is about $-\tau\,y\,T_k$: exactly the
shape of a constant repo difference of $-\tau y$.
\end{enumerate}
Two more shapes complete the picture: a contributor drop-off at pillar $T_k$ gives a step at $T_k$
(\cref{prop-nested}), and an FO repo mismark of $\Delta q$ gives $-\Delta q(T_k)\,T_k$, a smooth repo difference.
\end{proposition}
\begin{proof}
Apply \eqref{eq-delta-exact} to perturbed forwards or dividends.
(i) Each contributor's forward becomes $\bigl(S_0(1+\varepsilon)-D\bigr)e^{(r-q)T}$, so $\ln F^{\tot}_k$ moves by
$\ln\bigl(1+\varepsilon S_0/(S_0-D_k)\bigr)$; for the near dividend, $D^{\tot}_k$ moves by $\delta\tilde d$ at every
pillar.
(ii) \Cref{lem-allocation} gives $\delta D^{\tot}_k=a_{kp}\,\delta\ell_p$, zero before the ex-date and roughly constant
after.
(iii) At $T_k$, a fraction $f$ of the option panel includes the dividend, so $\ln F^{\tot}_k$ falls by
$f\tilde d_i/(S_0-D_k)$, while the allocation excludes it; before $T_k$, nobody includes it, and after $T_{k+1}$
everybody does. For the variance: the contributors form two groups, with weights $f$ and $1-f$, whose average
$\ln F_j$ differ by $\tilde d_i/(S_0-D_k)$.
(iv) $D^{\tot}_k-D_k=\sum_i\tilde d_i\bigl(\one_{\{t_i+\Delta\le T_k\}}-\one_{\{t_i\le T_k\}}\bigr)$, which equals
$-\tilde d_i$ exactly when $t_i\le T_k<t_i+\Delta$.
(v) Immediate. (vi) $\delta D^{\tot}_k=\tau D_k$; for a yield, $D_k=S_0(1-e^{-yT_k})\approx S_0\,y\,T_k$ and
$S_0-D_k\approx S_0$.
\end{proof}

\begin{table}[ht]
\centering
\caption{Signatures of the mechanisms. The last column gives the typical shape tested by the detector of
\cref{sec-detection}.}
\label{tab-signatures}
\small
\begin{tabular}{@{}>{\raggedright\arraybackslash}p{4.2cm}>{\raggedright\arraybackslash}p{4.3cm}
>{\raggedright\arraybackslash}p{3.6cm}l@{}}
\toprule
Mechanism & Carry difference $\delta$ & Repo difference & Typical shape\\
\midrule
Panel gap (\ref{prop-gap}) & smooth noise, increments at ex-dates & roughly flat noise & noise\\
Contributor drop-off (\ref{prop-nested}) & step at the drop-off & jump decaying like $1/T$ & step\\
(i) Spot mismatch & constant & hyperbola $-\varepsilon/T$ & level\\
(ii) Error on one period & step after the ex-date & jump decaying like $1/T$ & step\\
(iii) Ex-date disagreement & one pillar & spike & spike\\
(iv) Stale profile & pillars near the ex-dates & yearly zigzag & season\\
(v) Forward outlier & one pillar & spike & spike\\
(vi) Proportional bias & staircase; $\approx-\tau yT$ for a stream & steps; a level for an index & step\\
FO mismark & $-\Delta q(T)\,T$, smooth & smooth level or slope & none\\
\bottomrule
\end{tabular}
\end{table}

\begin{filrouge}
\Cref{fig-gallery} shows the running example's stock under each of the eight benchmark scenarios, with the same
contributor noise as in the clean market. In (a), the option contributors' spot is off by \vRunSpotEps: the Totem curve
dives like $1/T$ at the short end. In (b), the Totem 2027 level is off by \vRunBucketDelta: a jump at nine months that
decays like $1/T$. In (c), \vRunTimingShare{} of the option panel put the dividend going ex \vRunTimingDays~days after
the six-month expiry before it: a spike at six months. In (d), the allocation profile is shifted by
\vRunSeasonShift~days: the May dividends move before the six- and eighteen-month expiries, hence two negative spikes.
In (e), the consensus $\ln F$ at two years is off by \vRunFwdShift\,bp. In (f), the option contributors left after the
three-year drop-off overstate later dividends by \vRunPanelDelta: the distortion is smooth and stays inside the band.
In (g), it is the FO repo that is off by \vRunFoShift\,bp, Totem being intact. In (h), every Totem period is biased by
\vRunGrossBias{} (gross/net bias).
\end{filrouge}

\begin{figure}[p]
\centering\includegraphics[width=\linewidth]{figures/\lang/gallery.pdf}
\caption{Running example under the eight scenarios: FO repo, Totem-implied repo, band of $\pm2$ standard deviations of
consensus noise around the FO, and, in green, the Totem curve repaired by the detector of \cref{sec-detection}
(anomaly removed). In (f), the detector flags nothing; in (g), it flags no Totem anomaly but a possible FO mismark.
\Cref{tab-running} details the diagnoses.}
\label{fig-gallery}
\end{figure}

\subsection{What the data cannot tell}\label{sec-nonid}
\begin{theorem}[Non-identifiability of the split]\label{thm-nonid}
Assume European option prices depend on a contributor's views only through its forwards. For every contributor $j$
of the option panel, replace $(q_j, D_j)$ by $(q_j+u_j, D'_j)$ with
$u_j(T)\,T=\ln\bigl(S_0-D'_j(T)\bigr)-\ln\bigl(S_0-D_j(T)\bigr)$. Then every option price is unchanged, hence every
statistic of the option panel (consensus, dispersion, contributor count); the statistics of the dividend panel are
unchanged too. The implied repo $q^{\tot}$ is therefore unchanged, while the option panel's average repo $\qbar$ moves
by the average of the $u_j$. Consequently, the panel gap $q^{\tot}-\qbar$ is not a function of the observable data.
\end{theorem}
\begin{proof}
The transformation leaves $\ln F_j=(r-q_j)T+\ln(S_0-D_j)$ unchanged for every $j$ and every $T$: each contributor's
option prices, and every statistic derived from them, are unchanged. The dividend panel is untouched. $q^{\tot}$ is
computed from these statistics only, so it does not move, while $\qbar$ moves by the average of the $u_j$, which can
take any value.
\end{proof}

\begin{corollary}\label{cor-nonid}
No estimator built from Totem consensus data alone (curve smoothing, clustering, neural network) can remove the panel
gap. One can, however, compute its size (\cref{sec-detection}), recognise the structured mechanisms by their signature
(\cref{prop-signatures}), and do without the split between repo and dividends where the exposure does not need it
(\cref{sec-control}).
\end{corollary}

\begin{remark}[What resolves the indeterminacy]
Information outside the consensus pins down one of the two sides directly: listed dividend futures, total return
futures, which quote a funding spread, and stock borrow fees for the short end. Contributors common to both services
shrink the gap (\cref{cor-gap}), and American options carry a little information on the split through early exercise.
\end{remark}

\begin{retenir}
\begin{itemize}[leftmargin=12pt]
\item Implied repo $=$ option panel repo $+$ panel gap. The panel gap is normal noise of computable size: a few basis
points at one year, one to several tens at five years.
\item Real anomalies have recognisable shapes in carry: level (spot), step (dividend period, contributor drop-off),
spike (ex-date, forward outlier), season (stale profile).
\item The option panel's split between repo and dividends cannot be identified from Totem: no statistical method,
however sophisticated, will recover it.
\end{itemize}
\end{retenir}
